% Generated by Listas/make_statements.py from pratica-lecture-07.md, item 14. Do not edit.
\textbf{Source.} Tadelis (2013), \emph{Game Theory: An Introduction}, Exercise 12.6 ``Reap and Weep'', p.~268 {[}PDF p.~287{]}. Reproduced verbatim. \textbf{Answer in the source:} no (none in the book; Princeton UP's website has about half, not checked).



\begin{quote}

\textbf{Reap and Weep:} A farmer owns some land that he can farm to produce crops. Farming output depends on the talent of the farmer. The farmer knows his talent, but the rest of the world knows only that a farmer's talent is uniformly distributed: \(\theta \in [0, 1]\). The farmer's payoff from farming his land is equal to his talent \(\theta\). Before setting up his farm, the farmer approaches the local manufacturing plant and offers to work on the production line. The farmer can ask the plant owner for any wage \(w \ge 0\), and the owner can reject (\(R\)) or accept (\(A\)) the offer. If the owner rejects the offer then the farmer must return home and settle into his farming. If the owner accepts the offer then the farmer's payoff is the wage \(w\), while the owner's payoff is given by the net value \(\frac{3}{2}\theta - w\), and this is common knowledge.



\begin{enumerate}

\def\labelenumi{\alph{enumi}.}

\item

  Define the set of pure strategies for each player and find the pure-strategy Bayesian Nash equilibria of this game.

\item

  Averaging over the type of farmer, what are the possible levels of \emph{social surplus} (sum of expected payoffs of the farmer and the owner in their potential relationship) from the equilibrium you derived in (a)?

\item

  A local policy maker who is advocating for the increase of social surplus is proposing to cut water subsidies to the farmers, which would imply that a farmer of type \(\theta\) would get a payoff of \(\frac{1}{2}\theta\) from farming his land. This policy has no effect on the productivity of manufacturing. Using the criterion of social surplus, can you advocate for this policy maker using equilibrium analysis?

\end{enumerate}

\end{quote}

